\documentclass[aspectratio=169]{beamer} % 16:9; change to 11pt,aspectratio=43 if you prefer

% Theme (pick one: Madrid, Berlin, CambridgeUS, etc.)
\usetheme{Madrid}

% Basic packages
\usepackage[utf8]{inputenc}
\usepackage{amsmath, amssymb}
\usepackage[round]{natbib}
\bibliographystyle{apalike}
% Title info
\title[BST 245]{Power Estimation for Between Cluster Randomized Factorial Trials for Binary Outcomes}
\author{Phuc Vu}
\institute{HSPH}
\date{\today} % or a fixed date
\AtBeginSection[]{
  \begin{frame}{Outline}
    \tableofcontents[currentsection]
  \end{frame}
}
\begin{document}

\begin{frame}{Marginal ICC for Binary Outcomes}

\begin{itemize}
  \item For two individuals $i \neq j$ in the same cluster:
  \[
    \rho_{\text{marg}}
    = \mathrm{Corr}(Y_{i}, Y_{j})
    = \frac{\mathrm{Cov}(Y_i, Y_j)}
           {\sqrt{\mathrm{Var}(Y_i)\,\mathrm{Var}(Y_j)}} .
  \]
  \medskip
  
  \item For binary outcomes,
  \[
    \mathrm{Var}(Y_i) = p(1-p).
  \]
  
  \item Interpretation:
    \begin{itemize}
      \item Probability-scale measure of within-cluster similarity.
      \item Automatically shrinks as $p \to 0$ or $p \to 1$.
      \item Used in nearly all CRT power formulas.
    \end{itemize}
\end{itemize}

\end{frame}
\begin{frame}{Latent ICC in Logistic Mixed Models}

\textbf{Random-intercept GLMM:}
\[
  \text{logit}\{\Pr(Y_{ij}=1)\}
   = \beta_0 + \beta^\top X_{ij} + u_j,
  \qquad
  u_j \sim N(0,\tau^2).
\]

\begin{itemize}
  \item Equivalent latent-variable view:
  \[
    Z_{ij} = X_{ij}\beta + u_j + \varepsilon_{ij}, 
    \qquad
    \varepsilon_{ij} \sim \text{Logistic}(0,\pi^2/3).
  \]

  \item The \textbf{latent ICC} is:
  \[
    \rho_{\text{latent}}
    = \frac{\tau^2}{\tau^2 + \pi^2/3}.
  \]

  \item Natural for simulations that use GLMMs as the data-generating model.
\end{itemize}

\end{frame}
\begin{frame}{Working Model for Power Calculations}

\begin{itemize}
  \item For power estimation, we work with a
        \textbf{working linear model} at the cluster level:
    \[
      Z_j = X_j \beta + \varepsilon_j.
    \]

  \item We do not assume normal outcomes —
        this is a \textbf{variance-based approximation}
        used only for power calculations.

  \item If the error variances
        $\mathrm{Var}(\varepsilon_j) = w_j$ are known or approximated,
        then the covariance of $\hat\beta$ is given by
    \[
      \mathrm{Var}(\hat\beta)
      = (X^\top W^{-1} X)^{-1}.
    \]

  \item The key question becomes:
        \textbf{how should $Z_j$ and $W$ be defined for binary CRTs?}
\end{itemize}

\end{frame}

\begin{frame}{The Idea}

\begin{itemize}
  \item We work at the \textbf{cluster level} and focus on the
        \textbf{logit of the cluster mean outcome}.

  \item Approximate the logistic mixed model by a
        \textbf{linear model}:
    \[
      \text{logit}(\bar Y_j)
      \;\approx\;
      X_j \beta + \varepsilon_j,
    \]
    where $\varepsilon_j$ captures clustering, nonlinearity,
    and unequal cluster sizes.

  \item \textbf{Key reduction:}
        Power estimation depends entirely on approximating
        \[
          \mathrm{Var}\{\text{logit}(\bar Y_j)\}.
        \]

  \item Once this variance is known, power follows from
        standard GLS / Wald theory for linear models.
\end{itemize}

\end{frame}

\begin{frame}{Two Routes to Approximate the Variance}

\begin{itemize}
  \item From the Idea, power estimation reduces to approximating
    \[
      \mathrm{Var}\{\text{logit}(\bar Y_j)\}.
    \]

  \item We consider \textbf{two complementary approaches}:

  \begin{enumerate}
    \item \textbf{Analytic approximation}
      \begin{itemize}
        \item Uses marginal means, ICCs, and the delta method.
        \item Fast, interpretable, and suitable for design.
      \end{itemize}

    \item \textbf{Monte Carlo approximation}
      \begin{itemize}
        \item Simulates cluster means from the full GLMM.
        \item More computationally intensive, but closer to the true model.
      \end{itemize}
  \end{enumerate}

  \item Both approaches produce a working weight matrix $W$
        for GLS-based power calculations.
\end{itemize}

\end{frame}


\begin{frame}{Analytic Approximation}

\begin{enumerate}
  \item \textbf{Specify the factorial structure:}
        effect-coded $2^k$ design, full set of interactions for truth.
        
  \item \textbf{Fix coefficients $\beta$} (main effects + selected interactions).

  \item \textbf{Compute conditional probabilities:}
        \[
        p_{\text{cond}} = \text{logit}^{-1}(\eta), \qquad
        \eta = \beta_0 + X\beta.
        \]

  \item \textbf{Marginalize over random intercepts:}
        \[
          p_{\text{marg}} = 
          \mathbb{E}\left[ \text{logit}^{-1}(\eta + u) \right],
          \quad u \sim N(0,\tau^2).
        \]

  \item \textbf{Convert latent ICC to marginal ICC} for each cluster using Goldstein.

  \item \textbf{Approximate logit-scale covariance and compute power:}
        use CRT variance + delta method to form a diagonal weight matrix $W$,
        obtain the working covariance
        \[
          \Sigma = (X^\top W^{-1} X)^{-1},
        \]
        and apply standard GLMM/Wald theory to a contrast $c^\top\beta$
        (with small-sample $t$ df) to estimate power.
\end{enumerate}

\end{frame}


\begin{frame}{Variance Approximation for Binary Outcomes}

\begin{itemize}
  \item For each cluster $j$ with marginal mean $p_j$ and size $m_j$,
        the variance of an individual outcome is
    \[
      \mathrm{Var}(Y_{ij}) = p_j(1 - p_j),
      \qquad i = 1,\dots,m_j.
    \]

  \item Under a marginal ICC $\rho_{\text{marg}}$, the variance of the
        \textbf{cluster mean} is approximated by
    \[
      \mathrm{Var}(\bar Y_j)
      \approx
      \frac{p_j(1 - p_j)}{m_j}
      \Big[\,1 + (m_j - 1)\rho_{\text{marg}}\,\Big].
    \]

  \item Applying the \textbf{delta method} to the logit of the mean gives
    \[
      \mathrm{Var}\{\text{logit}(p_j)\}
      \approx
      \frac{\mathrm{Var}(\bar Y_j)}{p_j^2(1 - p_j)^2}.
    \]

  \item These per–cluster (or per–cell) logit variances define the diagonal
        working weight matrix
    \[
      W = \mathrm{diag}(w_1,\dots,w_J).
    \]
\end{itemize}

\end{frame}
\begin{frame}{The Final Piece: Linking Latent and Marginal ICC}

\begin{itemize}
  \item Power formulas for clustered binary data require
        \textbf{marginal ICC} $\rho_{\text{marg}}$.
  \item But simulations based on GLMMs specify \textbf{latent ICC}
        via the random intercept variance $\tau^2$.
  \medskip

  \item We therefore need a bridge between the two.

  \item Goldstein’s approximation:
    \[
      \rho_{\text{marg}}
      \;\approx\;
      \frac{\tau^2\,p(1-p)}
           {1 + \tau^2\,p(1-p)},
    \]
    where
    \[
      \rho_{\text{latent}}
      = \frac{\tau^2}{\tau^2 + \pi^2/3}.
    \]
\end{itemize}

\end{frame}
\begin{frame}{Unequal Cluster Sizes: Mean and SD Approximation}

\begin{itemize}

  \item Real CRTs often have \textbf{unequal cluster sizes}.
  
  \item Let cluster sizes follow some distribution with
    \[
      \bar m = \text{mean cluster size}, 
      \qquad
      s_m = \text{SD of cluster size}.
    \]

  \item Define the \textbf{coefficient of variation}:
    \[
      \text{CV} = \frac{s_m}{\bar m}.
    \]

  \item Standard design-effect adjustments replace $m_j$ by an
        \textbf{effective cluster size}:
    \[
      m^\star = \bar m(1 + \text{CV}^2).
    \]

  \item Our working variance for cluster means becomes:
    \[
      \mathrm{Var}(\bar Y_j)
        \;\approx\;
        \frac{p_j(1-p_j)}{\bar m}
        \Big[\,1 + (m^\star - 1)\rho_{\text{marg}}\,\Big].
    \]

  \item This matches common CRT practice (Eldridge \& Kerry; Hemming et al.).
\end{itemize}

\end{frame}
\begin{frame}{Monte Carlo Approximation}

\begin{itemize}
  \item Route number 2 is, we directly approximate
    \[
      \mathrm{Var}\{\text{logit}(\bar Y_j)\}
    \]
    under the full logistic mixed model.

  \item For each factorial cell $j$:
    \begin{enumerate}
      \item Draw a random intercept $u_j \sim N(0,\tau^2)$.
      \item Draw a cluster size $m_j$ from a specified distribution.
      \item Generate $Y_{ij} \sim \text{Bernoulli}(\text{logit}^{-1}(\eta_j + u_j))$.
      \item Compute $\bar Y_j$ and $\text{logit}(\bar Y_j)$.
    \end{enumerate}

  \item Repeating this many times yields an empirical estimate of
    \[
      w_j = \mathrm{Var}\{\text{logit}(\bar Y_j)\}.
    \]

  \item These variances define an alternative working weight matrix $W$
        for GLS-based power calculations.
\end{itemize}

\end{frame}

\begin{frame}{Power calculation given the Covariance matrix}

\begin{itemize}
  \item Expand the $2^k$ design to $K$ clusters using a 
        \textbf{balanced cluster allocation}.

  \item Working independence covariance:
    \[
      \Sigma 
      = (X^\top W^{-1} X)^{-1}.
    \]

  \item This gives standard errors for any contrast:
    \[
      \mathrm{SE}(\hat\theta)
      = \sqrt{c^\top \Sigma\, c}.
    \]

  \item For a main effect (e.g.\ component 1),  
        $c$ selects the corresponding coefficient.
\end{itemize}

\end{frame}
\begin{frame}{Small-Sample Correction and Power}

\begin{itemize}

  \item We test 
    \[
      H_0 : \theta = 0
      \quad\text{vs}\quad
      H_1 : \theta \neq 0.
    \]

  \item Wald statistic:
    \[
      t = \frac{\hat\theta}{\mathrm{SE}(\hat\theta)}.
    \]

  \item \textbf{Small-sample df correction:}
    \[
      \text{df} = K - p,
    \]
    where $p$ is the number of fixed-effect parameters
    in the fitted model (main + 2-way).

  \item Power:
    \[
      \text{Power}
      = 
      \Pr(|T_{\text{df},\lambda}| > t_{1-\alpha/2}),
      \qquad
      \lambda = \frac{\theta}{\mathrm{SE}}.
    \]

\end{itemize}

\end{frame}
\begin{frame}{Analytic vs.\ Monte Carlo Approximation}

\begin{itemize}
  \item \textbf{Analytic approximation}
    \begin{itemize}
      \item Fast and closed-form; well suited for design exploration.
      \item Relies on a delta-method approximation:
        \[
          \bar Y_j \;\approx\; \text{Normal}(p_j, \mathrm{Var}(\bar Y_j)).
        \]
      \item Accuracy improves as cluster size increases.
      \item May be less reliable when clusters are very small or probabilities
            are near 0 or 1.
    \end{itemize}

  \medskip

  \item \textbf{Monte Carlo approximation}
    \begin{itemize}
      \item Directly reflects the true GLMM data-generating process.
      \item Does not rely on normal approximations for $\bar Y_j$.
      \item Computationally more expensive; requires repeated simulation. (still take less than $1$s to run)
      \item Useful as a benchmark or when analytic assumptions are questionable.
    \end{itemize}

  \medskip

  \item \textbf{Practical recommendation:}
    \begin{itemize}
      \item Use the analytic approximation when average cluster sizes are moderate
            to large.
      \item Use the Monte Carlo approximation when cluster sizes are small
            or for sensitivity analysis.
    \end{itemize}
\end{itemize}

\end{frame}

\end{document}
